\chapter{Auswertung}
\label{ch:Auswertung}
% Liste der genutzer Formeln für die Fehlerrechnung
\section*{Fehlerrechnung}
Für die statistische Auswertung von $n$ Messwerten $x_i$ werden folgende Größen definiert \cite{errorSkript25}:
\begin{align}
    \bar{x} &= \frac{1}{n} \sum_{i=1}^{n} x_i \vphantom{\sqrt{\sum_i^n}^2} && \text{\textcolor{gray}{Arithmetisches Mittel}} \label{eq:arithmetisches_mittel} \\
    \sigma^2 &= \frac{1}{n-1} \sum_{i=1}^{n} (x_i - \bar{x})^2 \vphantom{\sqrt{\sum_i^n}^2} && \text{\textcolor{gray}{Varianz}} \label{eq:Varianz} \\
    \sigma &= \sqrt{\frac{1}{n-1} \sum_{i=1}^{n} (x_i - \bar{x})^2} \vphantom{\sqrt{\sum_i^n}^2} && \text{\textcolor{gray}{Standardabweichung}} \label{eq:standardabweichung} \\
    \Delta \bar{x} &= \frac{\sigma}{\sqrt{n}} = \sqrt{\frac{1}{n(n-1)} \sum_{i=1}^n(\bar x - x_i)^2} \vphantom{\sqrt{\sum_i^n}^2} && \text{\textcolor{gray}{Fehler des Mittelwerts}} \label{eq:fehler_mittelwert} \\
    \Delta f &= \sqrt{\left(\frac{\partial f}{\partial x} \Delta x\right)^2 + \left(\frac{\partial f}{\partial y} \Delta y\right)^2} \vphantom{\sqrt{\sum_i^n}^2} && \text{\textcolor{gray}{Gauß’sches Fehlerfortpflanzungsgesetz für $f(x,y)$}} \label{eq:gauss_fehlfortpflanzung} \\
    \Delta f &= \sqrt{(\Delta x)^2 + (\Delta y)^2} \vphantom{\sqrt{\sum_i^n}^2} && \text{\textcolor{gray}{Fehler für $f = x + y$}} \label{eq:fehler_summe} \\
    \Delta f &= |a| \Delta x \vphantom{\sqrt{\sum_i^n}^2} && \text{\textcolor{gray}{Fehler für $f = ax$}} \label{eq:fehler_proportional} \\
    \frac{\Delta f}{|f|} &= \sqrt{\left(\frac{\Delta x}{x}\right)^2 + \left(\frac{\Delta y}{y}\right)^2} \vphantom{\sqrt{\sum_i^n}^2} && \text{\textcolor{gray}{relativer Fehler für $f = xy$ oder $f = x/y$}} \label{eq:relativer_fehler} \\
    \sigma &= \frac{|a_{lit} - a_{gem}|}{\sqrt{\Delta a_{lit}^2 + \Delta a_{gem}^2}} \vphantom{\sqrt{\sum_i^n}^2} && \text{\textcolor{gray}{Berechnung der signifikanten Abweichung}} \label{eq:signifikante_abweichung}
\end{align}

\newpage

\section{Plateaubereich des Zählrohrs}
\label{sec:A1}
In dieser Sektion soll das Platau des Geiger-Müller-Zählrohres genauer untersucht werden. Die Daten der für diese Auswertung sind dem \hyperref[ch:Protokoll]{Messprotokoll} den Tabellen II.1, sowie III.1 zu entnehmen. 

Das bestimmte Pltau ist der \cabb{plots/Platau_1a} zu entnehmen. 
\img[1]{plots/Platau_1a}{Anzahl an Counts in Abhängigkeit der Betriebsspannung und visualisierung des Plataus. Eingezeichnet ist die Zählerstromspannung $U_0$ in der Mite des Plateaus.}

Zunächst soll die Differenz $D$ zwischen den Ereignissen $n_0$ bei $U_0 = (470 \pm 10)$V und $n_1$ bei $U_1 = (U_0 + 100)$V bestimmt werden. Die Ungenaugkeit der Differenz setzt sich aus den Ungenauigkeitden der Ereignisse wie folgt zusammen:
\eq{deltaDiff}{
    \Delta D = \sqrt{(\Delta n_0)^2 + (\Delta n_1)^2} \, .
}
Dabei kann sich zunutze gemacht werden, dasss die Messwerte (annährend) possionverteilt sind, dadurch lassen sich die Ungeauigkeiten der Häufigkeiten nach 
\eq{deltaNs}{
    \Delta n_i = \sqrt{n_i}
}
berechnen. Diese und folgende Ergebnisse sin der \ctab{A1_1} zu entnehmen. So soll auch das Signifikanzniveau $\zeta$ bestimmt werden, dieses kann via
\eq{sigNiv}{
    \zeta = \frac{D}[\Dleta D]
}
bestimmt werden. Die Ergebnisse sind der \ctab{A1_1} zu entnehmen, und zeigen, dass für die drei-, sowie die einminuten Messung die Steigung des Plateaus nicht signifikant war. Interessant ist, dass eine Verlängerung der Messzeit einen Trend zeigt, denn hier wird die Abweichung nochmal kleiner. Ob dies ein tatsächlicher Trend oder nur eine statistischer Zufall ist, wird in der \cch{Diskussion} diskutiert.
\com{Für die Diskussion; Eigentlich sollte die Standardabweichung für drei Minuten kleiner werden, da hier mehr Messungen und somit ein größerer prozentualer Fehler vorliegen. Dies ist also inkohärent mit der Theorie.}

Neben dem Signifikanzniveau kann auch die prozentuale Abweichung $\varrho$ bestimmt werden.
Die Formel für die prozentuale Abweichung wurde mit der \hyperref[eq:gauss_fehlfortpflanzung]{Gauß'schern Fehlerfortpflanzung (\ref*{eq:gauss_fehlfortpflanzung})} versehen und wird via
\eq{prozAbw}{
    \varrho = \frac{n_0}{n_1} \, , \quad
    \Delta \varrho = 
            \varrho \cdot \sqrt{\frac{\Delta {n_0}^{2}}{{n_0}^{2}} + \frac{\Delta {n_1}^{2}}{{n_1}^{2}}}
}
berechnet. Beide Werte liegen bei $\approx 1\%$.

\begin{table}[h!]
    \centering
    \label{tab:A1_1}
    \caption{Resultate der Berechnung zu den Eigenschaften des Plauteaus.}
    \begin{tabular}{l | C C C}
        \toprule
        Messzeit [s] & \text{Differenz } D & \text{Porzentuale Abweichung } n_0/n_1 [\%] & \text{Signifikanzniveau } \zeta [\sigma] \\
        \midrule
         60 & 90 \pm 100 & 1.018 \pm 0.020 & 0.89 \\
        180 & 40 \pm 180 & 0.997 \pm 0.011 & 0.22 \\
        \bottomrule
    \end{tabular}
\end{table}

\com{Hier ist noch etwas ausstehend. Es muss eine Messzeit $t$ bestimmt werden. Luca hat hier fancy Stuff gemacht, ist das pflicht oder doch nur Luca Nerd stuff? Und irgendwie Stuff um alle drei Sigma umgebungen, kp, muss du dir anschauen du noob.}

\section{Daten mit hoher mittlerer Ereigniszahl}
\label{sec:A2}
Diese Sektion wertet das Histogramm für einen Hohenmittelwert von Ereignissen aus. Dabei wurden hier 2500 Messungen durchgeführt (Dies ist im \hyperref[ch:Protokoll]{Messprotokoll} flasch notiert). Die Torzeit wurde auf 0.5s gesetzt. Es wurden die ersten 5 und die letzten 8 Messwerte heruasgeommen, da diese eine Häufigkeit von 0 hatten und dies die Regression erschwert/verfälscht, da so die Konvergenz gegen null nicht ordentlich dargestellt werden kann. 
Die Messdaten, sowie die Regressionen der Possion- und der Gaussverteilung sind der \cabb{plots/histo2b} zu entnehmen. Die Possionverteilung wird gestrichelt dargestellt, da diese nicht kontinuirlich ist. Die Fitparameter für beide Verteilungen sind ebenfalls dem Plot zu entnehmen. Zum Prüfen, wie gut diese Regressionen die Messdaten wieder spiegeln wird nach \csec{gutePlot} die Güte der Plots über das $\chi^2$ nach \ceq{chisquare} bestimmt. 
Dieses muss dann noch reduziert werden, es ergibt sich via \ceq{chisquare_red} das reduzierte Chi-Quadrat $\chi^2_\text{red}$. Zum reduzieren werden die Freiheitsgrade des jeweiligen Fits benötigt. Die Freiheitsgrade können via
\eq{DoF}{
    \text{DoF} = n_\text{mess} - n_\text{par}
}
berechnet werden, dabei beschreibt $n_\text{mess}$ die Mächtigkeit der Messdaten und $n_\text{par}$ die Anzahl der Fitparameter -- Für die Gaussverteilung gilt also $n_\text{par,g} = 3$ und für die Possionverteilung $n_\text{par,p} = 2$.
Der Zeilwert für einen guten Fit liegt bei rund um 1. Die Resultate sind in der \ctab{A2_1} eingetragen.
\begin{table}[h!]
    \centering
    \label{tab:A2_1}
    \caption{Bestimmung der Güte der Fits. Alle Werte auf zwei Dezimalstellen gerundet.}
    \begin{tabular}{l | C C C}
        \toprule
        Regression & \chi^2 & \chi^2_\text{red} & \text{Prozentuale Abweichung } [\%] \\
        \midrule
        Possion & 52.20 & 1.00 & 46.60 \\
        Gauss & 54.70 & 1.07 & 33.63 \\
        \bottomrule
    \end{tabular}
\end{table}
Zuletzt kann die Güte einer Regression noch über die Wahrscheinlichkeit eines Fits bestimmt werden. Diese Methode nimmt die bestimmten Regressionsparameter als wahr an und verteilt die Messwerte normal um den Mittelwert. Die Übereinstimmung der theoretischen und experimentellen Werte wird dann bestimmt. Auch diese Ergebnisse sind in der \ctab{A2_1} zu finden.
\img[1]{plots/histo2b}{Messdaten mit Regressionsfunktionen: Gauss- und Possionverteilungen bei 2500 Messungen.}

Die Ergebnisse sollen genau in der Diskussion analyisert werden.

\section{Daten mit kleinerer mittlerer Ereigniszahl}
\label{sec:A3}
In dieser Sektion werden dieselben Werte wie in \csec{A2} bestimmt, jedoch ist der Versuchsaufbau so verändert worden, dass hier der Mittelwert sehr viel geringer ausfällt. Die \cabb{plots/histo3b} zeigt das Histogramm und die Regressionen. 
\img[1]{plots/histo3b}{Histogramm von 6000 Messdaten mit Regressionsfunktionen: Gauss- und Possionverteilungen.}

Für diese Sektion wurde die Torzeit auf 0.1s gestellt und es wurden 6000 Messungen durchgeführt. Die Ergebnisse der Rechnung sind der \ctab{A3_1} zu entnehmen.

\begin{table}[h!]
    \centering
    \label{tab:A3_1}
    \caption{Bestimmung der Güte der Fits. Alle Werte auf zwei Dezimalstellen gerundet.}
    \begin{tabular}{l | C C C}
        \toprule
        Regression & \chi^2 & \chi^2_\text{red} & \text{Prozentuale Abweichung } [\%] \\
        \midrule
        Possion & 19.63 & 1.51 & 10.48 \\
        Gauss  & 128.15 & 10.68 & 0.0 \\
        \bottomrule
    \end{tabular}
\end{table}

